\documentclass[10pt,letterpaper]{article}
\usepackage{cornell-notes}

\title{Clause 33.05.10: Flag Type And Operations}
\author{}
\date{}
\setDocTitle{Clause 33.05.10: Flag Type And Operations}
\setDocSubtitle{C++ 2024}
\setDocOwner{C++ 2024 Cornell Notes}
\setDocDate{\today}
\setCornellCollection{C++ 2024 Cornell Notes}
\setCornellUnitType{Clause}
\setCornellUnitNumber{33.05.10}
\setCornellUnitTitle{Flag Type And Operations}
\hypersetup{pdftitle={Clause 33.05.10: Flag Type And Operations},pdfsubject={Cornell notes},pdfauthor={}}

\begin{document}
\maketitle
\begin{CornellOverviewBox}{Overview}
\textbf{Classification:} Standard library - Normative\\[2pt]
\textbf{Learning objective:} Understand the rules, terminology, and semantic consequences associated with Flag type and operations.
\end{CornellOverviewBox}\begin{CornellSummaryBox}{Summary}
{\sffamily\bfseries\color{Accent} Summary}\par\vspace{3pt}Section 33.5.10 focuses on Flag type and operations. Apply the specified interface together with its mandates, effects, result conditions, exception behavior, and complexity constraints. When applying it, preserve the distinction between mandatory requirements, recommendations, permissions, and behavior deliberately left undefined, unspecified, or implementation-defined.
\end{CornellSummaryBox}

\end{document}
